\section{Responsabilidade Civil e a Cadeia de Culpabilidade}
\label{sec:responsabilidade}

\subsection{Quem responde pela Iatrogenia Virtual?}

A adoção de gêmeos digitais introduz o fenômeno da responsabilidade compartilhada e difusa. Quando um implante perfura o canal mandibular ou uma reabsorção radicular severa ocorre apesar de um planejamento virtual "ótimo", a doutrina jurídica depara-se com múltiplos potenciais réus:
1. O \textbf{cirurgião-dentista}, que executou o ato ou validou o plano acriticamente;
2. A \textbf{empresa de software/IA}, cujos algoritmos calcularam trajetórias errôneas;
3. O \textbf{centro radiológico}, que forneceu um arquivo DICOM com artefatos suprimidos artificialmente que corromperam a segmentação.

No Brasil, o Código de Defesa do Consumidor (CDC) imputa responsabilidade subjetiva (dependente de culpa) ao profissional liberal, mas responsabilidade objetiva (independente de culpa, art. 12) ao fornecedor do software. Contudo, se a falha for decorrente da natureza opaca da rede neural profunda (\textit{Black Box}), provar o nexo de causalidade torna-se uma tarefa hercúlea para os tribunais \cite{springer2026realising}.

\destaque{A ausência de "Audit Trails" (trilhas de auditoria) imutáveis que detalhem o que o algoritmo recomendou \textit{versus} o que o dentista aprovou é a maior vulnerabilidade civil na atual prática da odontologia digital.}

\subsection{A Emergência de Seguros Cibernéticos Profissionais}

As apólices tradicionais de seguro de responsabilidade civil profissional foram redigidas para um paradigma analógico e frequentemente contêm cláusulas de exclusão veladas para iatrogenias causadas primariamente por falha sistêmica de software ou invasão cibernética. Um ataque de \textit{ransomware} que criptografe ou adultere um Gêmeo Digital armazenado localmente, levando à execução de um planejamento comprometido, não estaria coberto. O mercado segurador brasileiro precisará estruturar, urgentemente, apólices de \textit{Cyber-Medical Malpractice} para amparar os adotantes desta tecnologia.
